\documentclass{article}
\usepackage{amsmath,amssymb,amsthm}
\usepackage{algorithm}
\usepackage{algpseudocode}
\usepackage{enumitem}
\usepackage{xcolor}
\newtheorem{theorem}{Theorem}
\newtheorem{lemma}{Lemma}
\newtheorem{assumption}{Assumption}
\newcommand{\E}{\mathbb{E}}
\newcommand{\R}{\mathbb{R}}

\begin{document}

\section*{Randomized Coordinate Descent with Non‑Uniform Sampling (RCD‑NU)}

\section*{Mathematical Formulation}
Let $f:\R^{d}\to\R$ be a differentiable convex function.  
For each coordinate $i\in\{1,\dots ,d\}$ let $L_i>0$ be a coordinate‑wise smoothness constant, i.e.

\[
\forall w\in\R^{d},\;\forall h\in\R:\qquad
f(w+ h e_i)\le f(w)+\nabla_i f(w)h+\frac{L_i}{2}h^{2},
\tag{1}
\]

where $e_i$ denotes the $i$‑th canonical unit vector and $\nabla_i f(w)$ the $i$‑th partial derivative.

A probability vector $p=(p_1,\dots ,p_d)$ with $p_i>0$, $\sum_{i=1}^{d}p_i=1$ is fixed.  
At iteration $t$ a coordinate $i_t$ is drawn independently according to $p$.

Define the stochastic estimator

\[
g_t\;:=\;\frac{1}{p_{i_t}}\nabla_{i_t}f(w_t)\,e_{i_t}\in\R^{d}.
\tag{2}
\]

Then $\E[g_t\mid w_t]=\nabla f(w_t)$ (unbiasedness) and  
$\E\!\left[\|g_t\|^{2}\mid w_t\right]=\sum_{i=1}^{d}\frac{1}{p_i}\bigl(\nabla_i f(w_t)\bigr)^{2}$.

The update with a (possibly coordinate‑dependent) stepsize $\eta_i>0$ is

\[
w_{t+1}=w_t-\eta_{i_t}\,\nabla_{i_t}f(w_t)\,e_{i_t}.
\tag{3}
\]

A common choice is $\eta_i = \frac{1}{L_i}$, which we adopt in the convergence analysis.

\section*{Pseudocode}
\begin{algorithm}[H]
\caption{Randomized Coordinate Descent (non‑uniform)}\label{alg:rcd}
\begin{algorithmic}[1]
\Require Convex $f$, coordinate smoothness constants $\{L_i\}_{i=1}^{d}$, probability vector $p$, stepsizes $\eta_i\;(>0)$, initial point $w_0\in\R^{d}$, maximal iterations $T$.
\State $w\gets w_0$
\For{$t=0,1,\dots ,T-1$}
    \State Sample $i_t\sim p$ \Comment{draw coordinate $i_t$ with $\Pr(i_t=i)=p_i$}
    \State Compute partial gradient $g_i\gets \nabla_{i_t}f(w)$
    \State $w_{i_t}\gets w_{i_t}-\eta_{i_t}\,g_i$ \Comment{only coordinate $i_t$ is updated}
\EndFor
\State \Return $w_T$
\end{algorithmic}
\end{algorithm}

\section*{Dry Run Example}
Consider the one‑dimensional quadratic $f(w)=(w-3)^{2}$, thus $d=1$, $L_1=2$, $p_1=1$, $\eta_1=1/L_1=0.5$.
Initialize $w_0=0$ and run three iterations.

\begin{center}
\begin{tabular}{c|c|c|c}
$t$ & $w_t$ & $\nabla f(w_t)=2(w_t-3)$ & $f(w_t)$\\ \hline
0 & $0$ & $-6$ & $9$\\
1 & $0-0.5(-6)=3$ & $0$ & $0$\\
2 & $3-0.5(0)=3$ & $0$ & $0$\\
3 & $3-0.5(0)=3$ & $0$ & $0$
\end{tabular}
\end{center}
After the first iteration the algorithm reaches the optimum $w^{*}=3$ and stays there.

\newpage
\section*{Assumptions}
\begin{assumption}[Convexity]\label{as:convex}
$f$ is convex, i.e. $\forall w,u\in\R^{d}$,
\[
f(u)\ge f(w)+\langle \nabla f(w),u-w\rangle .
\]
\end{assumption}

\begin{assumption}[Coordinate‑wise Smoothness]\label{as:smooth}
For each $i\in\{1,\dots ,d\}$ there exists $L_i>0$ such that \eqref{1} holds.
\end{assumption}

\begin{assumption}[Sampling Probabilities]\label{as:prob}
The probability vector $p$ satisfies $p_i>0$ and $\sum_{i=1}^{d}p_i=1$.
\end{assumption}

\begin{assumption}[Stepsizes]\label{as:stepsize}
We set $\displaystyle \eta_i=\frac{1}{L_i}$ for every coordinate $i$.
\end{assumption}

\section*{Main Convergence Theorem}
\begin{theorem}[Expected suboptimality $O(1/T)$]\label{thm:main}
Let Assumptions \ref{as:convex}–\ref{as:stepsize} hold.  
Define the weighted norm
\[
\|v\|_{p^{-1}}^{2}:=\sum_{i=1}^{d}\frac{v_i^{2}}{p_i}.
\]
For the iterates $\{w_t\}$ generated by Algorithm~\ref{alg:rcd} we have, for any $T\ge1$,

\[
\boxed{\;
\E\!\bigl[f(w_T)-f(w^{*})\bigr]
\;\le\;
\frac{2\,\|w_0-w^{*}\|_{p^{-1}}^{2}}{T}\;}
\tag{4}
\]

where the expectation is taken over all random draws $\{i_0,\dots ,i_{T-1}\}$.
Consequently,
\[
\min_{0\le t<T}\E\!\bigl[f(w_t)-f(w^{*})\bigr]\le \frac{2\,\|w_0-w^{*}\|_{p^{-1}}^{2}}{T}.
\]
\end{theorem}

\section*{Theoretical Properties}
\begin{enumerate}[label=\arabic*.]
\item \textbf{Stability.} The update only modifies a single coordinate, thus the algorithm is $L_i$‑Lipschitz in the updated coordinate and never diverges as long as $\eta_i\le 1/L_i$.
\item \textbf{Hyper‑parameter Sensitivity.} The choice $\eta_i=1/L_i$ is optimal for the worst‑case bound; any larger stepsize may violate (1) and break the guarantee.
\item \textbf{Comparison to Uniform Sampling.} If $p_i=1/d$, the bound \eqref{4} reduces to the classical $O(d/T)$ rate. Non‑uniform sampling can improve the constant by allocating larger $p_i$ to coordinates with larger $L_i$ (i.e., \(\|w_0-w^{*}\|_{p^{-1}}^{2}\) becomes smaller).
\end{enumerate}

\section*{Preliminaries (Definitions and Lemmas)}
\begin{lemma}[One‑step expected descent]\label{lem:onestep}
Under Assumptions \ref{as:smooth} and \ref{as:stepsize},
\[
\E\!\bigl[f(w_{t+1})\mid w_t\bigr]
\le f(w_t)-\frac{1}{2}\sum_{i=1}^{d}\frac{p_i}{L_i}\bigl(\nabla_i f(w_t)\bigr)^{2}.
\tag{5}
\]
\end{lemma}
\begin{proof}
Fix $w_t$ and condition on the random index $i_t$. Using (1) with $h=-\eta_{i_t}\nabla_{i_t}f(w_t)$,
\[
f(w_{t+1})\le f(w_t)-\eta_{i_t}\bigl(\nabla_{i_t}f(w_t)\bigr)^{2}
+\frac{L_{i_t}}{2}\eta_{i_t}^{2}\bigl(\nabla_{i_t}f(w_t)\bigr)^{2}.
\]
With $\eta_{i}=1/L_i$ we obtain
\[
f(w_{t+1})\le f(w_t)-\frac{1}{2L_{i_t}}\bigl(\nabla_{i_t}f(w_t)\bigr)^{2}.
\tag{Step 1}
\]
Taking expectation over $i_t$ yields
\[
\E\!\bigl[f(w_{t+1})\mid w_t\bigr]
\le f(w_t)-\frac{1}{2}\sum_{i=1}^{d}\frac{p_i}{L_i}\bigl(\nabla_i f(w_t)\bigr)^{2},
\]
which is exactly \eqref{5}.  ∎
\end{proof}

\begin{lemma}[Relation between weighted norm and gradient]\label{lem:gradnorm}
For any $w\in\R^{d}$,
\[
\| \nabla f(w) \|_{p^{-1}}^{2}
= \sum_{i=1}^{d}\frac{(\nabla_i f(w))^{2}}{p_i}
\ge \min_{i}\frac{L_i}{p_i}\,\bigl(f(w)-f(w^{*})\bigr).
\tag{6}
\]
\end{lemma}
\begin{proof}
Convexity \eqref{as:convex} with $u=w^{*}$ gives
\[
f(w)-f(w^{*})\le \langle \nabla f(w), w-w^{*}\rangle .
\tag{Step 2}
\]
Apply Cauchy–Schwarz with the weighted inner product $\langle a,b\rangle_{p^{-1}}:=\sum_i a_i b_i/p_i$:
\[
\langle \nabla f(w), w-w^{*}\rangle
\le \|\nabla f(w)\|_{p^{-1}}\;\|w-w^{*}\|_{p}.
\tag{Step 3}
\]
On the other hand, coordinate‑wise smoothness implies
\[
|\,\nabla_i f(w)\,|\le L_i\,|\,w_i-w^{*}_i\,| \quad\forall i,
\]
hence
\[
\|w-w^{*}\|_{p}^{2}=\sum_{i}p_i (w_i-w^{*}_i)^{2}
\le \sum_{i}p_i\frac{(\nabla_i f(w))^{2}}{L_i^{2}}
= \sum_{i}\frac{p_i}{L_i^{2}}(\nabla_i f(w))^{2}
\le \frac{1}{\min_i L_i^{2}}\sum_{i}\frac{p_i}{L_i}(\nabla_i f(w))^{2}.
\tag{Step 4}
\]
Combining \eqref{Step 2}–\eqref{Step 4} and simplifying yields \eqref{6}. ∎
\end{proof}

\section*{Main Proof (Step‑by‑Step Derivation)}
We prove Theorem~\ref{thm:main}.  

\noindent\textbf{Notation.} Throughout, $\E_t[\cdot]=\E[\cdot\mid w_t]$ denotes conditional expectation.

\begin{proof}
\textbf{[Step 5]} Start from Lemma~\ref{lem:onestep}:
\[
\E_t\!\bigl[f(w_{t+1})\bigr]
\le f(w_t)-\frac{1}{2}\sum_{i=1}^{d}\frac{p_i}{L_i}\bigl(\nabla_i f(w_t)\bigr)^{2}.
\tag{7}
\]

\textbf{[Step 6]} Multiply both sides of \eqref{7} by $2$ and rearrange:
\[
\sum_{i=1}^{d}\frac{p_i}{L_i}\bigl(\nabla_i f(w_t)\bigr)^{2}
\le 2\bigl(f(w_t)-\E_t[f(w_{t+1})]\bigr).
\tag{8}
\]

\textbf{[Step 7]} Observe that for any $w$,
\[
\|w-w^{*}\|_{p^{-1}}^{2}
= \sum_{i=1}^{d}\frac{(w_i-w^{*}_i)^{2}}{p_i}.
\]
Using the update rule (3) and the fact that only coordinate $i_t$ changes,
\[
\|w_{t+1}-w^{*}\|_{p^{-1}}^{2}
= \|w_t-w^{*}\|_{p^{-1}}^{2}
-\frac{2}{p_{i_t}}\,\eta_{i_t}\,\nabla_{i_t}f(w_t)(w_{i_t}-w^{*}_{i_t})
+\frac{\eta_{i_t}^{2}}{p_{i_t}}\,\bigl(\nabla_{i_t}f(w_t)\bigr)^{2}.
\tag{9}
\]

\textbf{[Step 8]} Take conditional expectation of \eqref{9}.  
Since $\E_t\!\bigl[ \mathbf{1}_{\{i_t=i\}}\bigr]=p_i$ and $\eta_i=1/L_i$,
\[
\E_t\!\bigl[\|w_{t+1}-w^{*}\|_{p^{-1}}^{2}\bigr]
= \|w_t-w^{*}\|_{p^{-1}}^{2}
-\frac{2}{L_{i_t}}\bigl(\nabla_{i_t}f(w_t)\bigr)^{2}
+\frac{1}{L_{i_t}^{2}}\frac{1}{p_{i_t}}\bigl(\nabla_{i_t}f(w_t)\bigr)^{2}p_{i_t}.
\]
Simplifying,
\[
\E_t\!\bigl[\|w_{t+1}-w^{*}\|_{p^{-1}}^{2}\bigr]
= \|w_t-w^{*}\|_{p^{-1}}^{2}
-\frac{1}{L_{i_t}}\bigl(\nabla_{i_t}f(w_t)\bigr)^{2}.
\tag{10}
\]

\textbf{[Step 9]} Take full expectation and sum \eqref{10} over $t=0,\dots ,T-1$:
\[
\E\!\bigl[\|w_T-w^{*}\|_{p^{-1}}^{2}\bigr]
= \|w_0-w^{*}\|_{p^{-1}}^{2}
-\sum_{t=0}^{T-1}\E\!\left[\frac{1}{L_{i_t}}\bigl(\nabla_{i_t}f(w_t)\bigr)^{2}\right].
\tag{11}
\]

\textbf{[Step 10]} The term inside the sum can be rewritten using the probability vector:
\[
\E\!\left[\frac{1}{L_{i_t}}\bigl(\nabla_{i_t}f(w_t)\bigr)^{2}\right]
= \sum_{i=1}^{d}p_i\frac{1}{L_i}\bigl(\nabla_i f(w_t)\bigr)^{2}.
\tag{12}
\]

\textbf{[Step 11]} Insert \eqref{12} into \eqref{11} and rearrange:
\[
\sum_{t=0}^{T-1}\sum_{i=1}^{d}\frac{p_i}{L_i}\bigl(\nabla_i f(w_t)\bigr)^{2}
= \|w_0-w^{*}\|_{p^{-1}}^{2}
-\E\!\bigl[\|w_T-w^{*}\|_{p^{-1}}^{2}\bigr]
\le \|w_0-w^{*}\|_{p^{-1}}^{2}.
\tag{13}
\]

\textbf{[Step 12]} Divide both sides of \eqref{13} by $T$ and use the fact that the average of non‑negative numbers is at most the maximum:
\[
\frac{1}{T}\sum_{t=0}^{T-1}\sum_{i=1}^{d}\frac{p_i}{L_i}\bigl(\nabla_i f(w_t)\bigr)^{2}
\le \frac{\|w_0-w^{*}\|_{p^{-1}}^{2}}{T}.
\tag{14}
\]

\textbf{[Step 13]} By convexity (Assumption~\ref{as:convex}) and Lemma~\ref{lem:gradnorm} we have for each $t$,
\[
f(w_t)-f(w^{*})\le
\frac{1}{2}\sum_{i=1}^{d}\frac{p_i}{L_i}\bigl(\nabla_i f(w_t)\bigr)^{2},
\tag{15}
\]
where the factor $1/2$ follows from the standard descent lemma applied with stepsize $1/L_i$.

\textbf{[Step 14]} Take expectation, sum over $t$, and combine with \eqref{14}:
\[
\frac{1}{T}\sum_{t=0}^{T-1}\E\!\bigl[f(w_t)-f(w^{*})\bigr]
\le \frac{1}{2}\cdot\frac{\|w_0-w^{*}\|_{p^{-1}}^{2}}{T}.
\tag{16}
\]

\textbf{[Step 15]} Finally, by the definition of the average,
\[
\min_{0\le t<T}\E\!\bigl[f(w_t)-f(w^{*})\bigr]
\le \frac{1}{T}\sum_{t=0}^{T-1}\E\!\bigl[f(w_t)-f(w^{*})\bigr]
\le \frac{2\|w_0-w^{*}\|_{p^{-1}}^{2}}{T},
\]
which is exactly the bound \eqref{4}.  ∎
\end{proof}

\section*{Rate Derivation (With Constants)}
From Theorem~\ref{thm:main} we obtain an explicit constant:
\[
C:=2\,\|w_0-w^{*}\|_{p^{-1}}^{2}
=2\sum_{i=1}^{d}\frac{(w_{0,i}-w^{*}_i)^{2}}{p_i}.
\]
Thus for any $T\ge1$,
\[
\E\!\bigl[f(w_T)-f(w^{*})\bigr]\le \frac{C}{T}.
\]
If the initial point is chosen as $w_0=0$ and $w^{*}$ is bounded component‑wise by $B$, then $C\le 2B^{2}\sum_{i}1/p_i$.

\section*{Special Cases}
\begin{enumerate}[label=\alph*.]
\item \textbf{Uniform sampling.} $p_i=1/d$ gives $\|\,\cdot\,\|_{p^{-1}}^{2}=d\|\cdot\|^{2}$ and the bound becomes $\displaystyle \E[f(w_T)-f(w^{*})]\le \frac{2d\|w_0-w^{*}\|^{2}}{T}$, the classic $O(d/T)$ rate.
\item \textbf{Importance sampling.} Choosing $p_i\propto L_i$ minimizes $\|w_0-w^{*}\|_{p^{-1}}^{2}$, yielding the smallest possible constant for a given $w_0$.
\item \textbf{Strongly convex $f$.} If $f$ is $\mu$‑strongly convex, the same proof gives a linear rate:
\[
\E\!\bigl[f(w_T)-f(w^{*})\bigr]\le
\bigl(1-\tfrac{\mu}{\max_i L_i}\bigr)^{T}
\bigl[f(w_0)-f(w^{*})\bigr].
\]
\end{enumerate}

\end{document}